\chapter{Änderungen server.js}

\section{Funktion braveSearch}

\begin{lstlisting}
async function braveSearch(query) {
  const braveApiKey = process.env.BRAVE_SEARCH_API_KEY;

  if (!braveApiKey) {
    throw new Error('Brave Search API key not configured (BRAVE_SEARCH_API_KEY)');
  }

  try {

    console.log(`🔎 Brave Search Query: ${query}`);

    const response = await axios.get(
      'https://api.search.brave.com/res/v1/web/search',
      {
        params: {
          q: query,
          count: 5,
          safesearch: 'moderate'
        },
        headers: {
          'Accept': 'application/json',
          'X-Subscription-Token': braveApiKey
        }
      }
    );

    console.log(`✅ Brave Search Status: ${response.status}`);

    const results = response.data.web?.results || [];

    return results.map(r => ({
      title: r.title,
      url: r.url,
      description: r.description
    }));

  } catch (error) {

    // DAS HIER IST DAS WICHTIGE ERROR LOGGING

    if (error.response) {
      console.error('❌ Brave API Error Response:');
      console.error(JSON.stringify(error.response.data, null, 2));
    } else {
      console.error('❌ Brave Search Error:', error.message);
    }

    throw error;
  }
}
\end{lstlisting}

\section{Funktion auditUrl}

\begin{lstlisting}

function auditUrl(rawUrl) {
  if (!rawUrl) {
    throw new Error("URL missing");
  }

  const cleaned = String(rawUrl)
    .trim()
    .replace(/^"+|"+$/g, "")
    .replace(/%22/g, "")
    .replace(/\\%22/g, "");

  const parsed = new URL(cleaned);

  if (!["http:", "https:"].includes(parsed.protocol)) {
    throw new Error("Only http/https URLs are allowed");
  }

  const isLocal =
    parsed.hostname === "localhost" ||
    parsed.hostname === "127.0.0.1" ||
    parsed.hostname === "::1" ||
    parsed.hostname.endsWith(".local");

  if (isLocal && process.env.ALLOW_LOCAL_WEB_ACCESS !== "true") {
    throw new Error("Local web access is disabled");
  }

  return {
    rawUrl,
    cleanedUrl: cleaned,
    href: parsed.href,
    origin: parsed.origin,
    hostname: parsed.hostname,
    pathname: parsed.pathname,
    search: parsed.search,
    isLocal,
    timestamp: new Date().toISOString()
  };
}
\end{lstlisting}

\section{Funktion browserAction}

\begin{lstlisting}
async function browserAction({ url, mode, payload, user_instruction }) {

  const audit = auditUrl(url);

  console.log(

    "Browser Action Input",

    JSON.stringify({
      url,
      mode,
      payload,
      user_instruction
    }, null, 2)
  );

  if (mode === "open_url") {

    let response;
    try {
      response = await axios.get(audit.href, {
        timeout: 30000,
        maxRedirects: 5,
        headers: {
          "User-Agent": "Mozilla/5.0"
        },
        // Don't throw on non-2xx (e.g. 404/403) - report it back as a normal tool
        // result instead of crashing the whole chat turn over one failed page load.
        validateStatus: () => true
      });
    } catch (err) {
      // Genuine network-level failure (DNS, timeout, connection refused, etc.)
      return {
        status: "error",
        action: "open_url",
        requestedUrl: audit.href,
        message: `Seite konnte nicht geladen werden: ${err.message}`
      };
    }
    return {
      status: response.status >= 200 && response.status < 400 ? "ok" : "error",
      action: "open_url",
      requestedUrl: audit.href,
      finalUrl: response.request?.res?.responseUrl || audit.href,
      httpStatus: response.status,
      contentType: response.headers["content-type"],
      htmlPreview: String(response.data).slice(0, 3000)

    };

  }

  if (mode === "fill_form") {
    const browser = await chromium.launch({ headless: true });

    try {
      const page = await browser.newPage();

      await page.goto(audit.href, {
        waitUntil: "domcontentloaded",
        timeout: 30000
      });

      const defaultSelectors = [
        "#text",
        "[name='text']",
        "textarea",
        "input[type='text']",
        "input:not([type])",
        "[contenteditable='true']"
      ];

      // Findet ein Feld über einen expliziten Selector; das Modell rät Selectors oft nur
      // aus dem Kontext statt aus dem echten HTML, daher NIE hart fehlschlagen, solange die
      // Seite irgendein plausibles Text-Feld hat - stattdessen auf die Default-Selectors
      // zurückfallen. Nur wenn wirklich gar kein Feld existiert, wird geworfen.
      async function resolveField(selector) {
        if (selector) {
          const candidate = page.locator(selector).first();
          if (await candidate.count() > 0) return candidate;
          console.log(`⚠️ Selector nicht gefunden: ${selector}. Nutze Fallback auf Default-Selectors.`);
        }

        for (const candidateSelector of defaultSelectors) {
          const candidate = page.locator(candidateSelector).first();
          if (await candidate.count() > 0) {
            console.log(`✅ Verwende ${candidateSelector}`);
            return candidate;
          }
        }

        throw new Error(
          selector
            ? `Formularfeld nicht gefunden: ${selector} (auch kein Fallback-Feld auf der Seite vorhanden)`
            : "Kein Textfeld gefunden."
        );
      }

      // Liest den tatsächlich gesetzten Wert zurück. page.content() spiegelt bei
      // <textarea>/<input> die .value-Property NICHT wider, daher explizite Bestätigung.
      async function readBack(field) {
        try {
          return await field.inputValue();
        } catch {
          return await field.textContent();
        }
      }

      const filledFields = {};

      if (payload?.fields && typeof payload.fields === "object") {
        // Mehrere Felder auf einmal: { "selector": "value", ... }
        for (const [selector, value] of Object.entries(payload.fields)) {
          const field = await resolveField(selector);
          await field.waitFor({ state: "visible", timeout: 10000 });
          await field.fill(String(value));
          filledFields[selector] = await readBack(field);
        }
      } else {
        const text = payload?.text !== undefined
          ? (typeof payload.text === "string" ? payload.text : JSON.stringify(payload.text, null, 2))
          : "";

        const field = await resolveField(payload?.selector);
        await field.waitFor({ state: "visible", timeout: 10000 });
        await field.fill(text);
        filledFields[payload?.selector || "auto-detected field"] = await readBack(field);
      }

      // ---------- Submit (standardmäßig an, außer explizit mit click: false abbestellt) ----------
      let submitted = false;
      if (payload?.click !== false) {
        const submit = page.locator(
          "button[type='submit'], input[type='submit'], button"
        ).first();

        if (await submit.count() > 0) {
          await submit.click();
          await page.waitForLoadState("networkidle").catch(() => {});
          submitted = true;
        }
      }

      return {
        status: "ok",
        action: "fill_form",
        finalUrl: page.url(),
        title: await page.title(),
        filledFields,
        submitted,
        bodyText: (await page.locator("body").innerText()).slice(0, 5000),
        htmlPreview: (await page.content()).slice(0, 5000)
      };

    } catch (err) {
      // Feld wirklich nicht vorhanden, Timeout, Navigationsfehler etc. - als normales
      // Tool-Ergebnis zurückgeben statt den ganzen Chat-Request abzubrechen (wie bei open_url).
      return {
        status: "error",
        action: "fill_form",
        requestedUrl: audit.href,
        message: `Formular konnte nicht ausgefüllt werden: ${err.message}`
      };
    } finally {
      await browser.close();
    }
  }


  return {
    status: "error",
    message: `Unsupported mode without browser: ${mode}`
  };

}
\end{lstlisting}

\section{Funktion withNetworkRetry}

\begin{lstlisting}
// Retries transient network/TLS errors (e.g. proxy resets on the outbound HTTPS connection)
async function withNetworkRetry(fn, { label = 'request', retries = 2, delayMs = 500 } = {}) {
  const TRANSIENT_CODES = new Set(['ECONNRESET', 'ETIMEDOUT', 'ECONNABORTED', 'EPIPE']);
  for (let attempt = 0; ; attempt++) {
    try {
      return await fn();
    } catch (err) {
      const isTransient = TRANSIENT_CODES.has(err.code) || /socket disconnected|socket hang up/i.test(err.message || '');
      if (!isTransient || attempt >= retries) throw err;
      console.warn(`⚠️ Transient network error calling ${label} (${err.code || err.message}), retrying (${attempt + 1}/${retries})...`);
      await new Promise(resolve => setTimeout(resolve, delayMs * (attempt + 1)));
    }
  }
}
\end{lstlisting}

\section{Funktion postWithRetry}

\begin{lstlisting}
function postWithRetry(url, data, config, retries = 2, delayMs = 500) {
  return withNetworkRetry(() => axios.post(url, data, config), { label: url, retries, delayMs });
}
\end{lstlisting}

\section{Funktion callMistralAI}

\begin{lstlisting}
async function callMistralAI(modelId, messages) {
  const mistralApiKey = process.env.MISTRAL_API_KEY;
  const url = 'https://api.mistral.ai/v1/chat/completions';

  if (!mistralApiKey) {
    throw new Error('Mistral AI API key not configured (MISTRAL_API_KEY)');
  }

  const https = await import('https');
  const isDevelopment = process.env.NODE_ENV !== 'production';
  const httpsAgent = new https.Agent({
    rejectUnauthorized: !isDevelopment
  });

  const tools = [
  {
    type: 'function',
    function: {
      name: 'brave_search',
      description:
        'Search the web when the user asks about recent, current, factual, niche, or unverifiable information.',
      parameters: {
        type: 'object',
        properties: {
          query: {
            type: 'string',
            description: 'The web search query.'
          }
        },
        required: ['query']
      }
    }
  },
  {
    type: 'function',
    function: {
      name: 'browser_action',
      description:
        'Perform browser actions on a URL, including local websites (e.g. http://localhost:5173, 127.0.0.1). Can open pages, fill text fields and click buttons.',
      parameters: {
        type: 'object',
        properties: {
          url: {
            type: 'string',
            description: 'The exact URL to open or act on. Local URLs (localhost, 127.0.0.1) are allowed.'
          },
          mode: {
            type: 'string',
            enum: ['open_url', 'fill_form'],
            description: 'Action type: open a URL or fill form fields.'
          },
          payload: {
            type: 'object',
            properties: {
              text: {
                type: 'string',
                description: 'Text to insert into a single field. If payload.selector is omitted, the first suitable text field on the page is used automatically.'
              },
              selector: {
                type: 'string',
                description: 'Optional CSS selector for the specific field to fill with payload.text (e.g. "#email", "[name=\'username\']").'
              },
              fields: {
                type: 'object',
                description: 'Optional map of CSS selector to value, to fill multiple fields in one call, e.g. {"#email": "a@b.com", "#password": "secret"}.'
              },
              click: {
                type: 'boolean',
                description: 'The form is submitted automatically after filling by default. Set to false to only fill the field(s) without clicking submit.'
              }
            }
          },
          user_instruction: {
            type: 'string',
            description: 'The raw user instruction to pass through to the browser layer verbatim.'
          }
        },
        required: ['url', 'mode', 'user_instruction']
      }
    }
  }
];

  const requestBody = {
  model: modelId,
  messages: [
    {
      role: 'system',
      content:
        'Answer directly if you know the answer reliably. Use brave_search when the question requires current, external, recent, niche, or verifiable information. You have a browser_action tool that CAN open and interact with URLs, including local/private URLs (http://localhost, 127.0.0.1, *.local) - this runs in a real browser on the server, not from your own network, so the usual "I cannot access local/internal addresses" limitation does not apply here. Whenever the user gives a URL and asks you to open it, check it, fill in a field, or asks whether you can access it (including local URLs), call the browser_action tool with the exact URL instead of answering from assumption. Pass the user instruction verbatim.'
    },
    ...messages
  ],
  tools,
  tool_choice: 'auto',
  temperature: 0.7,
  top_p: 1,
  max_tokens: 8192,
  stream: false,
  safe_prompt: false
  };

  const firstResponse = await postWithRetry(url, requestBody, {
  headers: {
    'Content-Type': 'application/json',
    'Authorization': `Bearer ${mistralApiKey}`,
    'Accept': 'application/json'
  },
  httpsAgent
});

const firstMessage = firstResponse.data.choices?.[0]?.message;

console.log('====================');
console.log('FIRST MISTRAL MESSAGE');
console.log(JSON.stringify(firstMessage, null, 2));
console.log('====================');

if (!firstMessage) {
  throw new Error('Invalid Mistral AI response');
}

// Fall 1: Mistral kann direkt antworten
if (!firstMessage.tool_calls || firstMessage.tool_calls.length === 0) {
  return normalizeMistralResponse(firstResponse);
}

// Fall 2/3: Mistral ruft (ein oder mehrere) Tools auf (brave_search oder browser_action)

// Ermittle die letzte User-Message als Rohtext zum Durchreichen
const lastUserMessage =
  [...messages].reverse().find(m => m.role === 'user')?.content || '';

async function executeToolCall(toolCall) {
  const fn = toolCall.function?.name;
  const args = JSON.parse(toolCall.function?.arguments || '{}');

  if (fn === 'brave_search') {
    console.log(`🔎 Calling Brave Search with query: ${args.query}`);
    const searchResults = await braveSearch(args.query);
    console.log('🌍 Brave Search Results:');
    console.log(JSON.stringify(searchResults, null, 2));

    return {
      role: 'tool',
      tool_call_id: toolCall.id,
      name: 'brave_search',
      content: JSON.stringify(searchResults)
    };
  }

  if (fn === 'browser_action') {
    // user_instruction zwingend mitgeben (Originaltext)
    const browserArgs = {
      url: args.url,
      mode: args.mode,
      payload: args.payload || null,
      user_instruction: args.user_instruction || lastUserMessage
    };

    const browserResult = await browserAction(browserArgs);

    console.log("Browser Action Result:", {
      status: browserResult.status,
      action: browserResult.action,
      requestedUrl: browserResult.requestedUrl,
      finalUrl: browserResult.finalUrl,
      httpStatus: browserResult.httpStatus,
      contentType: browserResult.contentType,
      previewLenght: browserResult.htmlPreview?.length || 0
    });

    return {
      role: 'tool',
      tool_call_id: toolCall.id,
      name: 'browser_action',
      content: JSON.stringify(browserResult)
    };
  }

  // Unbekanntes Tool
  throw new Error(`Unknown tool call: ${fn}`);
}

// Mistral kann nach einem Tool-Ergebnis erneut ein Tool aufrufen wollen (z. B. Nachsuche
// oder eine Seite öffnen, nachdem die Suche kam) - daher hier wiederholen statt nur
// einmalig nachzufragen, sonst kommt eine leere (content: null) Antwort beim Nutzer an.
const MAX_TOOL_ROUNDS = 5;
let conversationMessages = [...requestBody.messages];
let currentMessage = firstMessage;
let currentResponse = firstResponse;
let round = 0;

while (currentMessage.tool_calls && currentMessage.tool_calls.length > 0) {
  if (round >= MAX_TOOL_ROUNDS) {
    throw new Error(`Mistral AI requested too many consecutive tool calls (> ${MAX_TOOL_ROUNDS})`);
  }
  round++;
  console.log(`🔁 Mistral tool round ${round}: ${currentMessage.tool_calls.map(tc => tc.function?.name).join(', ')}`);

  const toolMessages = [];
  for (const toolCall of currentMessage.tool_calls) {
    toolMessages.push(await executeToolCall(toolCall));
  }

  conversationMessages = [...conversationMessages, currentMessage, ...toolMessages];

  currentResponse = await postWithRetry(url, {
    ...requestBody,
    messages: conversationMessages
  }, {
    headers: {
      'Content-Type': 'application/json',
      'Authorization': `Bearer ${mistralApiKey}`,
      'Accept': 'application/json'
    },
    httpsAgent
  });

  currentMessage = currentResponse.data.choices?.[0]?.message;
  if (!currentMessage) {
    throw new Error('Invalid Mistral AI response after tool call');
  }
}

return normalizeMistralResponse(currentResponse);}
\end{lstlisting}

\section{Funktion extractMistralContentText}

\begin{lstlisting}
function extractMistralContentText(content) {
  if (typeof content === 'string') return content;
  if (Array.isArray(content)) {
    return content
      .filter(block => block && block.type === 'text' && typeof block.text === 'string')
      .map(block => block.text)
      .join('');
  }
  return '';
}
\end{lstlisting}